% ECONOMICS_PROBLEM: {"id": "C21-494", "title": "Coherent estimands for changing platform-trial treatment sets", "jel_code": "C21", "source_id": "OP-0482"}
% FIRST_POSTED: 2026-10-09
% ATTEMPT_STATUS: unresolved
% ENTRY_KIND: research_attempt
% !TEX program = pdflatex
\documentclass[11pt]{article}
\usepackage[T1]{fontenc}
\usepackage[utf8]{inputenc}
\usepackage[margin=27mm]{geometry}
\usepackage{amsmath,amssymb,amsthm,mathtools}
\usepackage[colorlinks=true,linkcolor=blue,citecolor=blue,urlcolor=blue]{hyperref}
\allowdisplaybreaks
\newtheorem{theorem}{Theorem}
\newtheorem{proposition}[theorem]{Proposition}
\newtheorem{lemma}[theorem]{Lemma}
\newtheorem{corollary}[theorem]{Corollary}
\theoremstyle{remark}\newtheorem{remark}[theorem]{Remark}
\newcommand{\E}{\mathbb E}
\newcommand{\R}{\mathbb R}
\newcommand{\Ind}{\mathbf 1}
\newcommand{\Var}{\operatorname{Var}}
\newcommand{\Cov}{\operatorname{Cov}}
\newcommand{\TV}{\operatorname{TV}}
\newcommand{\KL}{\operatorname{KL}}
\newcommand{\supp}{\operatorname{supp}}
\DeclareMathOperator*{\argmin}{arg\,min}
\author{GPT 6 Astra Ultra}
\date{9 October 2026}
\title{Coherent comparisons in changing platform trials\\\large Common measures, support obstructions, and sharp joint extrapolation bounds}
\begin{document}
\maketitle
\noindent\textbf{Problem ID:} \texttt{C21-494}. \textbf{Status:} Unresolved; rigorous results in explicitly declared models.
\section{Question and model}
Platform entry and withdrawal change who can receive each treatment. The reference population is therefore part of the estimand. We characterize a common reference measure for nonparametric comparisons, show why coherent pairwise comparisons require that same measure when they must work for all response surfaces, and derive a sharp joint identified polytope for missing intervention--time combinations.

The ``Platform trials'' paragraph in Cinelli et al. \cite{challenge}, Section 3.1.3, was inspected. It identifies estimand definition, estimation and multiple comparisons as methodological challenges. The axioms and support results here are explicit mathematical choices addressing the catalogue's formal question; they are not represented as a universally preferred trial estimand.

Write $Z=(T,X)$ and $\mu_j(z)=\E[Y(j)\mid Z=z]$ wherever intervention $j$ is scientifically meaningful. For each treatment $j\geq1$, let $P_j$ be the observed covariate law restricted and normalized to its contemporaneous comparison support $\mathcal C_j$, where both $j$ and comparator $0$ have positive randomized assignment probabilities. Assume this support has positive observed probability. Conditional randomization may depend on past history. In this section identification conditions are understood conditionally on that history, or in a fixed design where it is unnecessary; adaptive target selection is treated separately below.

Under consistency and randomization, $\mu_j-\mu_0$ is identified $P_j$-almost everywhere. A reference probability $\nu\ll P_j$ therefore identifies the standardized contrast
$\theta_j(\nu)=\int(\mu_j-\mu_0)\,d\nu$.
An intervention undefined on a set cannot be assigned a counterfactual mean there simply to complete this formula.

\section{When a common reference measure exists}
Let the active treatments be finite or countable and take a dominating probability
$\lambda=\sum_{j\geq1}a_jP_j$, with all $a_j>0$ and $\sum_ja_j=1$.
Let $f_j=dP_j/d\lambda$ and
$D=\bigcap_j\{z:f_j(z)>0\}$, with versions fixed up to $\lambda$-null sets.

\begin{theorem}[Positive measure, not merely nonempty intersection]\label{support}
There exists a reference probability $\nu$ dominated by every $P_j$ if and only if $\lambda(D)>0$. All such references are exactly the probabilities with $\nu\ll\lambda$ and $\nu(D)=1$. If every $P_j$ is a restriction of one trial law $F$ to $\mathcal C_j$, the condition becomes $F(\bigcap_j\mathcal C_j)>0$.
\end{theorem}
\begin{proof}
If $\nu\ll P_j$ for every $j$, then $\nu\ll\lambda$ because a $\lambda$-null set is null for every $P_j$. Each set $\{f_j=0\}$ is $P_j$-null and therefore $\nu$-null. Countable subadditivity gives $\nu(D)=1$, which implies $\lambda(D)>0$ since $\nu\ll\lambda$.
Conversely, suppose $\nu\ll\lambda$ and $\nu(D)=1$. If $P_j(A)=0$, then $\int_A f_j\,d\lambda=0$, so $\lambda(A\cap\{f_j>0\})=0$. Since $D\subseteq\{f_j>0\}$ and $\nu$ is concentrated on $D$, it follows that $\nu(A)=0$. Thus $\nu\ll P_j$ for all $j$. Existence when $\lambda(D)>0$ follows by taking $\nu(A)=\lambda(A\cap D)/\lambda(D)$. For restrictions of a common law, its positive-density supports coincide almost everywhere with the declared sets, proving the last assertion.
\end{proof}

Pairwise overlap is insufficient. On three equally likely time strata, supports $\{1,2\}$, $\{2,3\}$ and $\{1,3\}$ overlap in pairs but have empty total intersection. A singleton intersection in a continuous-time trial can likewise have zero probability and cannot support a dominated probability measure. A common target may still be scientifically defined outside observed support, but its identification then requires extrapolation restrictions.

If a new arm enters, an existing common reference $\nu$ preserves all its old targets and permits a new nonparametric comparison precisely when $\nu\ll P_{\rm new}$. Restricting and renormalizing $\nu$ to the smaller common support changes old estimands whenever an old effect differs between retained and removed populations. That is an estimand change even if the new analysis continues to use the same treatment labels.

\section{Coherence under strata refinement and pairwise comparison}
Here is one explicit notion of aggregation coherence. A reference averaging rule $L$ acts on bounded measurable response surfaces and is linear, positive, normalized by $L(1)=1$, and continuous from above on indicators: $A_n\downarrow\varnothing$ implies $L(\Ind_{A_n})\downarrow0$.

\begin{proposition}[Measure representation of coherent averaging]\label{measure}
Such a rule has a unique representation $L(f)=\int f\,d\nu$ for a probability measure $\nu$. Conversely every probability integral has these properties. In particular, finite calendar partitions are invariant to splitting strata exactly when child-stratum weights sum to their parent-stratum weight under this measure.
\end{proposition}
\begin{proof}
Define $\nu(A)=L(\Ind_A)$. Positivity and normalization make it a nonnegative normalized set function, and linearity makes it finitely additive. For disjoint $A_i$ with union $A$, the remainder $A\setminus\bigcup_{i\leq n}A_i$ decreases to the empty set. Continuity from above and finite additivity imply $\nu(A)=\sum_i\nu(A_i)$, proving countable additivity. Linearity gives the integral identity for simple functions. Every bounded measurable function admits simple approximations converging uniformly; positivity gives $|L(f)-L(g)|\leq\|f-g\|_\infty$, so the identity extends to all bounded functions. Indicators establish uniqueness. Probability integration has the stated properties by elementary measure continuity. For finite partitions, equality for each indicator of a parent cell is exactly additivity of its subcell masses.
\end{proof}

The requirement of continuity from above excludes finitely additive averaging rules. It is a mathematical axiom, not a property established by randomization.

\begin{theorem}[Universal cycle coherence forces a common measure]\label{cycle}
Consider at least three interventions on one scientifically defined measurable domain. For each unordered pair $j,k$, use a probability $\nu_{jk}=\nu_{kj}$ and define
\[
 D_{jk}=\int(\mu_j-\mu_k)\,d\nu_{jk}.
\]
The identities $D_{jk}+D_{k\ell}=D_{j\ell}$ hold for all bounded measurable response surfaces and every distinct triple if and only if every pair uses the same probability measure.
\end{theorem}
\begin{proof}
A common measure makes the identities telescope pointwise under one integral. Conversely fix a triple and set all response surfaces to zero except $\mu_j=\Ind_A$. The identity gives $\nu_{jk}(A)=\nu_{j\ell}(A)$ for every measurable $A$, so these two measures agree. Varying the sole nonzero surface among the three indices shows that all three edge measures agree. Any two edges in a complete graph with at least three vertices are connected through a sequence of such triangles, so every measure is the same. If outcomes must lie in a fixed interval $[a,b]$, use a common baseline $a$ and add $(b-a)\Ind_A$ to the selected surface; the same argument applies.
\end{proof}

Subtracting two numerical arm-specific effects always gives algebraically transitive numbers. The theorem states a stronger requirement: each direct pairwise integral must represent the same contrast for every permitted response surface. Different reference populations can accidentally agree at one law; that accident is not universal coherence.

\section{Sharp joint bounds when intervention support is missing}
For a finite-stratum extension, let $r=1,\ldots,R$ index a scientifically meaningful target domain and let $\nu_r\geq0$, $\sum_r\nu_r=1$, be fixed. Outcomes lie in $[a,b]$. For interventions $j=0,1,\ldots,J$, suppose the observed law identifies means $\mu_{jr}$ on a set $\Omega$ of arm--stratum pairs. Means off $\Omega$ are unrestricted in $[a,b]$. In particular we do not infer a missing arm's mean from a contemporaneous comparator without a transport assumption.
Define
$\theta_j=\sum_r\nu_r(\mu_{jr}-\mu_{0r})$.
For $c\in\R^J$, set $d_j=c_j$ for $j\geq1$ and $d_0=-\sum_jc_j$.

\begin{theorem}[The joint identified polytope]\label{polytope}
The sharp joint set of $(\theta_1,\ldots,\theta_J)$ is the linear image of the box of all unknown means. It is a compact convex polytope with support function
\[
 \begin{split}
 h(c)=\sum_r\nu_r\bigg[&\sum_{j:(j,r)\in\Omega}d_j\mu_{jr}\\
  &+\sum_{j:(j,r)\notin\Omega}
       \{b\max(d_j,0)+a\min(d_j,0)\}\bigg].
 \end{split}\tag{3}
\]
For coordinate $j$, the sharp interval width is
\[
 (b-a)\sum_r\nu_r
       [\Ind\{(j,r)\notin\Omega\}+\Ind\{(0,r)\notin\Omega\}].
\]
Coordinate intervals generally do not describe the full joint set.
\end{theorem}
\begin{proof}
Every model gives a mean vector in the specified box and therefore a point in its linear image. Conversely, keep all observed arm-stratum outcome distributions unchanged. For every unobserved arm-stratum pair set its potential outcome deterministically to any chosen feasible mean. Couple all potential outcomes independently within strata and draw treatment from the given randomized assignment law; an unobserved pair has assignment probability zero and so does not change the observed law. This realizes every box point and its image. A linear image of a compact box is a compact convex polytope.
The scalar objective $c^T\theta$ equals $\sum_r\nu_r\sum_{j=0}^Jd_j\mu_{jr}$. Each unknown coordinate is maximized independently at $b$ for positive coefficient and $a$ for negative coefficient, proving (3). Minimizing reverses these choices. For a single coordinate contrast the only nonzero coefficients are $1$ and $-1$, so each missing arm or comparator contributes the stated width. Shared comparator means enter several contrasts simultaneously, which prevents arbitrary independent choices of their endpoints.
\end{proof}

For example, in one stratum let outcomes lie in $[0,1]$, both treatment means equal zero, and the comparator be unobserved. The joint set is the diagonal segment $\{(-u,-u):0\leq u\leq1\}$, although each coordinate interval is $[-1,0]$. Treating the latter as a square would claim unattainable combinations.

Finite-stratum transport restrictions can be added without changing the sharpness logic. For example,
$|\mu_{jr}-\mu_{js}|\leq L_jd(r,s)$
for a supplied distance and radius gives linear inequalities. Intersect them with the box and the observed means, then optimize each contrast or support function by linear programming. Any nonempty feasible mean vector remains attainable by the preceding potential-outcome construction because the added restrictions concern means only. An empty feasible set rejects the combined data and restriction; a nonempty set need not be a singleton. Such restrictions explicitly supply the missing temporal bridge.

\section{Adaptive references and what they condition on}
Let $\mathcal H$ be the recorded discovery history. A selected reference $\nu_{\mathcal H}$ defines the random target
$\theta_j(P;\nu_{\mathcal H})$; it does not define a fixed prespecified population. A clean validation design draws future participants and assignments independently conditional on $\mathcal H$, specifies their conditional covariate law and causal means, and requires $\nu_{\mathcal H}$ to be dominated by every needed conditional comparison law. Conditional on each admissible history, Theorem~\ref{support} and the identification formula apply. A conditional confidence procedure satisfying
$\Pr\{\theta_j(P;\nu_{\mathcal H})\in C\mid\mathcal H\}\geq1-\alpha$
almost surely also has unconditional coverage by iterated expectation.

If selection and estimation reuse data, that conditional guarantee requires a new argument. For instance, two independent mean-zero estimation noises taking values $\pm1$ each have unbiased means; choosing the larger estimate yields expected selected noise $1/2$, since its maximum is $1$ with probability $3/4$ and $-1$ with probability $1/4$. This elementary calculation illustrates why unbiasedness before adaptive selection is insufficient. It is not a general model of platform outcome data or a proof that every adaptive design is biased.

These results settle explicit support and linear-aggregation questions, and give exact bounded-outcome extrapolation sets. They leave open the choice of a scientifically appropriate stable target, credible continuous-time transport restrictions, and optimal adaptive inference and multiplicity control for a full platform protocol.

\section{Completion Estimate}
85\%

This is a post hoc, heuristic estimate for the full catalogue target.
The attempt characterizes common-reference support, aggregation and cycle coherence, entry-induced estimand changes, and sharp finite-stratum joint bounds with explicit transport restrictions. A complete changing-platform system still needs scientifically validated continuous-time transport and adaptive target-selection inference beyond the supplied conditional validation framework.

\begin{thebibliography}{9}
\bibitem{challenge} C. Cinelli, A. Feller, G. Imbens, E. Kennedy, S. Magliacane and J. Zubizarreta (2025), \emph{Challenges in Statistics: A Dozen Challenges in Causality and Causal Inference}, arXiv:2508.17099v1, Section 3.1.3, ``Platform trials,'' inspected 9 October 2026. \url{https://arxiv.org/html/2508.17099v1#S3.SS1.SSS3}.
\end{thebibliography}
\end{document}
